\documentclass[11pt]{letter}

\usepackage{geometry}
\geometry{margin=1in}
\usepackage{setspace}
\usepackage{hyperref}

\signature{Decory Edwards}
\address{Decory Edwards \\ Johns Hopkins University \\ Department of Economics}

\begin{document}

\begin{letter}{Hiring Committee \\ Hudson River Trading}

\opening{Dear Hiring Committee,}

I am writing to apply for the Algorithm Developer / Quant Researcher position for 2026 PhD graduates at Hudson River Trading. I am currently a PhD candidate in Economics at Johns Hopkins University. My training combines mathematical reasoning, computational modeling, and data driven empirical analysis. I am motivated by problems that require precise logic, careful implementation, and constant validation against real data. HRT’s emphasis on rigorous research, fast feedback, and collaborative problem solving strongly aligns with how I approach research and coding.

My research agenda is at the intersection of wealth inequality and macroeconomics. In my job market paper, I build and estimate a structural heterogeneous agent model with returns that differ across households. The core contribution is to show how heterogeneity in returns and income risk jointly shape the wealth distribution and consumption behavior. Relative to closely related work, my framework performs better at matching the lower and middle portions of the wealth distribution. This difference matters quantitatively. Models that focus primarily on returns heterogeneity can fit the top of the wealth distribution closely but tend to generate too much wealth among the bottom half of households. In my framework, income uncertainty generates a precautionary saving motive that produces genuinely low wealth households. This leads to a realistic distribution of marginal propensities to consume and aggregate consumption responses that align with empirical estimates.

A key feature of my research is the heavy use of computation to solve and simulate models that cannot be handled analytically. I write code to solve dynamic programming problems, simulate large populations, and evaluate model fit using moments from survey data such as the SCF and HRS. This work requires careful numerical reasoning, performance awareness, and constant debugging. Beyond the job market paper, my broader research portfolio reflects a consistent approach to problem solving. I work with large panel datasets, construct new variables from raw microdata, and design empirical strategies that isolate economically meaningful variation. My empirical work relies on econometric reasoning and causal inference.

I am particularly interested in Hudson River Trading because of its research culture and the way engineers and researchers collaborate closely. The role’s focus on developing and refining algorithms, analyzing large scale data, and iterating based on performance feedback fits well with my training. My PhD has emphasized writing code that is transparent, testable, and interpretable. I value environments where ideas are challenged directly and improved through discussion and evidence. HRT’s flat structure and emphasis on intellectual honesty are especially appealing.

My economics training has provided a strong foundation in probability, optimization, and statistical inference. I am comfortable reasoning from first principles, proving results when needed, and translating theory into working code. I have extensive experience using Python for numerical methods, simulation, and data analysis. I also regularly work with version control and reproducible workflows. These tools are integral to my research process and would be central to my contributions at HRT.

I am also enthusiastic about the prospect of living and working in one of HRT’s core offices. I value in person collaboration and the fast learning that comes from working closely with other researchers and engineers. I am comfortable relocating and would welcome the opportunity to build my career in a city with a strong research and technology community. I see this role as a long term investment in developing as a quantitative researcher in a demanding and rewarding environment.

I would welcome the opportunity to contribute my skills and perspective to Hudson River Trading. I am excited by the challenge of applying rigorous quantitative reasoning to real time decision making. Thank you for your time and consideration.

\closing{Sincerely,}

\end{letter}
\end{document}
